\newcommand{\xliftto}[1]{\mathrel{\smash{\stackrel{\raisebox{3.4pt}{\tiny $#1\:$}}{\smash{\liftto}}}}}
\newcommand{\xleftliftto}[1]{\mathrel{\smash{\stackrel{\raisebox{3.4pt}{\tiny $\:#1$}}{\smash{\leftliftto}}}}}
\newcommand{\liftto}{\rightrightarrows}
\newcommand{\iliftto}{\xliftto{\mathbf{i}}}

\chapter{Probabilistic Term Rewrite Systems}
\chapterauthor{Jan-Christoph Kassing}

\begin{abstract}
  This chapter describes the syntax and semantics for first-order \emph{Probabilistic Term Rewrite Systems} (PTRSs).
  PTRSs extend ordinary TRSs by allowing rewrite rules with finitely many probabilistic outcomes.
  In the Termination Competition, the categories \emph{PTRS Standard} and \emph{PTRS Innermost}
  use the same underlying input format and differ only in whether arbitrary rewrite steps or only innermost rewrite steps are considered.
\end{abstract}

\section{Syntax}

\subsection{Grammar}

\begin{tabular}{lll}
ptrs & ::= & \texttt{(format PTRS)} fun* prule* \\
fun & ::= & \texttt{(fun} IDENTIFIER NAT\texttt{)} \\
term & ::= & IDENTIFIER $|$ \texttt{(}IDENTIFIER term+\texttt{)} \\
prule & ::= & \texttt{(prule} term dist\texttt{)} \\
dist & ::= & \texttt{((} outcome+\texttt{))} \\
outcome & ::= & \texttt{(}term prob?\texttt{)} \\
prob & ::= & \texttt{:prob} NAT \\
\end{tabular}

\bigskip
Additionally, we have NAT $\in \mathbb{N}$, and IDENTIFIER are intended to be valid SMT-LIB symbols.
As in the TRS format, quoted and unquoted identifiers should not be mixed, and reserved words such as
\texttt{fun}, \texttt{prule}, and \texttt{format} must not be used as identifiers.
Function symbols are declared by \texttt{(fun f n)}, where \texttt{n} is the arity of \texttt{f}.
Identifiers that are not declared as function symbols are variables.

Each probabilistic rule contains a finite list of outcomes.
The optional annotation \texttt{:prob q} assigns the positive weight $q$ to the corresponding outcome.
If it is omitted, the default weight is $1$.
The actual branch probabilities are obtained by normalizing all weights of the rule.
For example, the outcomes \texttt{((a :prob 2) (b :prob 1))} represent the distribution
$\{\tfrac{2}{3}:\mathsf{a},\tfrac{1}{3}:\mathsf{b}\}$.

\textbf{Example.}
The PTRS $\mathcal{P} = \{\mathsf{f} \to \{\tfrac{1}{2}:\mathsf{f},\; \tfrac{1}{2}:\mathsf{stop}\}\}$
has the following input syntax:
\begin{center}
\begin{tabular}{l}
\texttt{(format PTRS)} \\
\texttt{(fun f 0)} \\
\texttt{(fun stop 0)} \\
\texttt{(prule f ((f :prob 1) (stop :prob 1)))}
\end{tabular}
\end{center}

\section{Semantics}

We write $\mathcal{T}(\F,\V)$ for the set of all terms over a signature $\F = \biguplus_{k \in \N} \F_k$ of function symbols and a set $\V$ of variables.
A function symbol $f \in \F_k$ has arity $k$.
As usual, substitutions, positions, subterms, and contexts are defined exactly as for ordinary TRSs.

A \emph{finite multi-distribution} on a set $A$ is a multiset
$\mu = \{p_1:a_1,\ldots,p_m:a_m\}$ where every $0 \leq p_i < 1$,
every $a_i \in A$, and $\sum_{i=1}^m p_i = 1$.
The set of all finite multi-distributions on $A$ is denoted $\mathit{FDist}(A)$.

A \emph{probabilistic rewrite rule} is a pair $\ell \to \mu$, where $\ell \in \mathcal{T}(\F,\V)$ and
$\mu$ is a finite multi-distribution over $\mathcal{T}(\F,\V)$.
A \emph{probabilistic term rewrite system} (PTRS) is a finite set $\Rules$ of probabilistic rewrite rules.

The concrete input syntax stores positive weights instead of normalized probabilities.
Thus, a declaration
\[
  \texttt{(prule }\ell\texttt{ ((}r_1\texttt{ :prob }w_1\texttt{) \ldots (}r_k\texttt{ :prob }w_k\texttt{)))}
\]
denotes the rule
\[
  \ell \to \Big\{\frac{w_1}{W}:r_1,\ldots,\frac{w_k}{W}:r_k\Big\},
  \qquad\text{where } W = \sum_{i=1}^{k} w_i.
\]

The PTRS $\Rules$ induces a relation
${\to_\Rules} \subseteq \mathcal{T}(\F,\V) \times \mathit{FDist}(\mathcal{T}(\F,\V))$.
We define $s \to_\Rules \mu$ if there exist a position $p \in \mathit{Pos}(s)$,
a rule $\ell \to \{q_1:r_1,\ldots,q_k:r_k\} \in \Rules$, and a substitution $\sigma$
such that $s|_p = \ell\sigma$ and
\[
  \mu = \{q_1:s[r_1\sigma]_p,\ldots,q_k:s[r_k\sigma]_p\}.
\]
Thus, one rewrite step replaces a redex by a finite probability distribution over successor terms.
A term $u$ is a \emph{redex} w.r.t.\ $\Rules$ if there exist a rule $\ell \to \mu \in \Rules$ and a substitution $\sigma$ with $u = \ell\sigma$.
A term is in \emph{normal form} w.r.t.\ $\Rules$ if it contains no redex as a subterm.
An \emph{innermost probabilistic rewrite step}, written $s \ito_\Rules \mu$, is a rewrite step $s \to_\Rules \mu$ at a position $p$
such that all proper subterms of the redex $s|_p$ are in normal form w.r.t.\ $\Rules$.

\subsection{Lifted Rewriting and Probabilistic Termination}

To describe termination in the probabilistic setting, rewrite steps are lifted from single terms to
multi-distributions over terms.

Let ${\to} \subseteq \mathcal{T}(\F,\V) \times \mathit{FDist}(\mathcal{T}(\F,\V))$ be a probabilistic relation.
For $0 < p \leq 1$ and $\mu = \{q_1:t_1,\ldots,q_m:t_m\}$, let $p \cdot \mu = \{p\cdot q_1:t_1, \ldots, p\cdot q_m:t_m\}$.
The \emph{lifting} of $\to$, written $\liftto$, is the smallest relation on multi-distributions satisfying the following clauses:
\begin{itemize}[itemsep=0pt]
  \item If $t$ is in normal form w.r.t.\ $\to$, then $\{1:t\} \liftto \{1:t\}$.
  \item If $t \to \mu$, then $\{1:t\} \liftto \mu$.
  \item If $\mu_1 \liftto \nu_1,\ldots,\mu_k \liftto \nu_k$ and $p_1,\ldots,p_k > 0$ with $\sum_{j=1}^{k} p_j = 1$, then
  \[
    \bigcup_{j=1}^{k} p_j \cdot \mu_j \liftto \bigcup_{j=1}^{k} p_j \cdot \nu_j.
  \]
\end{itemize}
For a PTRS $\Rules$, we write $\liftto_\Rules$ and $\iliftto_\Rules$for the lifting of $\to_\Rules$ and $\ito_\Rules$, respectively.

For a multi-distribution $\mu$, let
\[
  |\mu|_{\Rules} = \sum_{(p:t) \in \mu,\; t \in \mathit{NF}_{\Rules}} p,
\]
that is, $|\mu|_{\Rules}$ is the total probability mass of normal forms.
Let $(\mu_n)_{n \in \mathbb{N}}$ be an infinite $\liftto_\Rules$-rewrite sequence,
that is, $\mu_n \liftto_\Rules \mu_{n+1}$ for all $n \in \mathbb{N}$.
The PTRS $\Rules$ is \emph{almost-surely terminating} (AST) if for every such sequence with
$\mu_0 = \{1:t\}$ for some term $t$, we have $\lim_{n \to \infty} |\mu_n|_{\Rules} = 1$.
Intuitively, every infinite rewrite process reaches a normal form with probability~$1$.

For an infinite $\liftto_\Rules$-rewrite sequence $\vec{\mu} = (\mu_n)_{n \in \mathbb{N}}$, its
\emph{expected derivation length} is defined by
\[
  \mathit{edl}(\vec{\mu}) = \sum_{n=0}^{\infty} (1 - |\mu_n|_{\Rules}).
\]
For a start term $t$, its \emph{expected derivation height} is
\[
  \mathit{edh}_{\Rules}(t) = \sup\{ \mathit{edl}(\vec{\mu}) \mid \vec{\mu}
  \text{ is a } \liftto_\Rules\text{-rewrite sequence with } \mu_0 = \{1:t\} \}.
\]
The PTRS $\Rules$ is \emph{strongly almost-surely terminating} (SAST) if
$\mathit{edh}_{\Rules}(t) < \infty$
for every start term $t$.
Equivalently, for every term $t$ there is a finite bound $C_t$ such that every lifted rewrite sequence
starting from $\{1:t\}$ has expected derivation length at most $C_t$.

The notions of \emph{innermost AST} and \emph{innermost SAST} are obtained by replacing
$\to_\Rules$ and $\liftto_\Rules$ with $\ito_\Rules$ and $\iliftto_\Rules$.

\section{Competition Categories}

The PTRS formalism is used by the following two categories.

\subsection{PTRS Standard}

The PTRS Standard category is concerned with the question whether every rewrite sequence
eventually stops with probability $1$ and whether every start term has a finite upper bound
on the expected runtime of all rewrite sequences starting from it.
In other words, the category asks for proofs of AST or of the stronger property SAST.

\begin{formalproblem}
  {PTRS Standard}
  {PTRS $\Rules$}
  {Is $\Rules$ AST or SAST?}
  {\texttt{WORST\_CASE(}$f$\texttt{,}$g$\texttt{)}, where}
  \[f,g \in \{\texttt{MAYBE} \sqsubseteq \texttt{AST} \sqsubseteq \texttt{SAST} \sqsubseteq \texttt{INF}\},\]
  and where $f$ is the lower bound and $g$ is the upper bound.
\end{formalproblem}

\textbf{Motivation.}
This category captures probabilistic termination without fixing an evaluation strategy.
It is the natural analogue of unrestricted rewriting for probabilistic programs,
where each rewrite step may branch into several successor states with different probabilities.

\textbf{Example.}
The PTRS $\mathcal{P} = \{\mathsf{g}(x) \to \{\tfrac{1}{2}:\mathsf{g}(\mathsf{g}(x)),\; \tfrac{1}{2}:x\}\}$
has the following input syntax:
\begin{center}
\begin{tabular}{l}
\texttt{(format PTRS)} \\
\texttt{(fun g 1)} \\
\texttt{(prule (g x) (((g (g x)) :prob 1) (x :prob 1)))}
\end{tabular}
\end{center}
This system induces a symmetric random walk on the number of outermost $\mathsf{g}$ symbols.
It is AST, since divergence has probability $0$, but it is not SAST because the expected derivation length is unbounded.
Hence, the expected output can be \texttt{WORST\_CASE(AST,AST)}.

The PTRS $\mathcal{P} = \{\mathsf{f}(x) \to \{\tfrac{1}{3}:\mathsf{f}(\mathsf{f}(x)),\; \tfrac{2}{3}:x\}\}$
has the following input syntax:
\begin{center}
\begin{tabular}{l}
\texttt{(format PTRS)} \\
\texttt{(fun f 1)} \\
\texttt{(prule (f x) (((f (f x)) :prob 1) (x :prob 2)))}
\end{tabular}
\end{center}
This random walk is biased towards decreasing the number of outermost $\mathsf{f}$ symbols.
Hence, every start term has finite expected derivation height, so the system is SAST.
Thus, the expected output can be \texttt{WORST\_CASE(SAST,SAST)}.

The PTRS $\mathcal{P} = \{\mathsf{h}(x) \to \{\tfrac{2}{3}:\mathsf{h}(\mathsf{h}(x)),\; \tfrac{1}{3}:x\}\}$
has the following input syntax:
\begin{center}
\begin{tabular}{l}
\texttt{(format PTRS)} \\
\texttt{(fun h 1)} \\
\texttt{(prule (h x) (((h (h x)) :prob 2) (x :prob 1)))}
\end{tabular}
\end{center}
This random walk is biased towards larger terms and is not AST.
Therefore, neither AST nor SAST can be certified, so the expected output should be \texttt{WORST\_CASE(INF,INF)}.

\subsection{PTRS Innermost}

The PTRS Innermost category restricts rewriting to innermost redexes.
Thus, it asks whether every innermost rewrite sequence eventually stops with probability $1$
and whether every start term has a finite upper bound on the expected runtime of all innermost rewrite sequences.

\begin{formalproblem}
  {PTRS Innermost}
  {PTRS $\Rules$}
  {Is $\Rules$ innermost AST or innermost SAST?}
  {\texttt{WORST\_CASE(}$f$\texttt{,}$g$\texttt{)}, where}
  \[f,g \in \{\texttt{MAYBE} \sqsubseteq \texttt{AST} \sqsubseteq \texttt{SAST} \sqsubseteq \texttt{INF}\},\]
  and where $f$ is the lower bound and $g$ is the upper bound.
\end{formalproblem}

\textbf{Motivation.}
This category studies probabilistic rewriting under a call-by-value style strategy.
Innermost rewriting is particularly relevant for probabilistic programming languages and interpreters
that evaluate arguments before applying a rule at the root.

\textbf{Example.}
The PTRS $\mathcal{P} = \{\mathsf{g}(x) \to \{\tfrac{1}{2}:\mathsf{g}(\mathsf{g}(x)),\; \tfrac{1}{2}:x\}\}$
has the following input syntax:
\begin{center}
\begin{tabular}{l}
\texttt{(format PTRS)} \\
\texttt{(fun g 1)} \\
\texttt{(prule (g x) (((g (g x)) :prob 1) (x :prob 1)))}
\end{tabular}
\end{center}
Since rewriting only happens at the root, full and innermost rewriting coincide.
Thus, this symmetric random walk is innermost AST, but not innermost SAST.
Hence, the expected output can be \texttt{WORST\_CASE(AST,AST)}.

The PTRS $\mathcal{P} = \{\mathsf{f}(x) \to \{\tfrac{1}{3}:\mathsf{f}(\mathsf{f}(x)),\; \tfrac{2}{3}:x\}\}$
has the following input syntax:
\begin{center}
\begin{tabular}{l}
\texttt{(format PTRS)} \\
\texttt{(fun f 1)} \\
\texttt{(prule (f x) (((f (f x)) :prob 1) (x :prob 2)))}
\end{tabular}
\end{center}
Again, all rewrite steps are root steps, so innermost and full rewriting agree.
Since the walk is biased towards decreasing the number of outermost $\mathsf{f}$ symbols, it is innermost SAST.
Thus, the expected output can be \texttt{WORST\_CASE(SAST,SAST)}.

The PTRS $\mathcal{P} = \{\mathsf{h}(x) \to \{\tfrac{2}{3}:\mathsf{h}(\mathsf{h}(x)),\; \tfrac{1}{3}:x\}\}$
has the following input syntax:
\begin{center}
\begin{tabular}{l}
\texttt{(format PTRS)} \\
\texttt{(fun h 1)} \\
\texttt{(prule (h x) (((h (h x)) :prob 2) (x :prob 1)))}
\end{tabular}
\end{center}
Here innermost rewriting also coincides with full rewriting.
This random walk is biased towards larger terms and is therefore not innermost AST.
Hence, the expected output should be \texttt{WORST\_CASE(INF,INF)}.